\chapter{I modelli}
\label{ch:modelli}

I sei modelli formano una catena: ciascuno estende il precedente con un elemento. Questo capitolo
li presenta in ordine, dando per ognuno le variabili introdotte, la funzione obiettivo e i vincoli.
Il significato dettagliato dei singoli termini dell'obiettivo è nel capitolo~\ref{ch:obiettivo}.

\section{Modello base --- \code{cvxpy}}
\label{sec:m1}

\noindent\textbf{Sorgente:} \code{approach\_cvxpy.py} $\rightarrow$ \code{constrained\_highs}\\
\textbf{Always-on:} stimati come baseline e sottratti; il solver vede il residuo\\
\textbf{Livelli di potenza:} 1 (acceso/spento)

\subsection*{Funzione obiettivo}

\begin{modelbox}{Obiettivo del modello base}
\begin{equation}
  \min \;\;
  \underbrace{\sum_{t \in \Vset} \bigl( \ep_t + \en_t \bigr)}_{\text{errore di ricostruzione}}
  \;+\;
  \underbrace{\sum_{i} \sum_{t \in \Vset} \lambda_w \, p_i \, f_i(t) \, x_{i,t}}_{\text{fascia oraria}}
  \;+\;
  \underbrace{\sum_{i} \sum_{t \in \Vset} \lambda_a \, p_i \, u_{i,t}}_{\text{accensioni}}
\end{equation}
\end{modelbox}

\subsection*{Vincoli}

\paragraph{Ricostruzione del segnale.} Un'uguaglianza per ogni slot valido, che lega la somma
delle potenze stimate al valore misurato attraverso le due variabili di errore:
\begin{equation}
  \sum_{i} p_i \, x_{i,t} \;-\; \ep_t \;+\; \en_t \;=\; y_t
  \qquad \forall t \in \Vset .
  \label{eq:ricostruzione-base}
\end{equation}

\paragraph{Coerenza delle transizioni.} Le variabili $u$ e $\dw$ devono corrispondere ai reali
cambi di stato di $x$. Il vincolo fondamentale è un bilancio di flusso: la variazione di stato fra
due slot consecutivi è pari alla differenza fra accensione e spegnimento.
\begin{align}
  u_{i,0} &= x_{i,0} && \forall i \label{eq:init-u}\\
  \dw_{i,0} &= 0 && \forall i \label{eq:init-dw}\\
  x_{i,t} - x_{i,t-1} - u_{i,t} + \dw_{i,t} &= 0 && \forall i,\; \forall t \ge 1 \label{eq:transizione}\\
  u_{i,t} + \dw_{i,t} &\le 1 && \forall i,\; \forall t \ge 1 \label{eq:mutex}
\end{align}
La~\eqref{eq:init-u} tratta il primo slot del giorno come un'accensione se il device risulta
acceso, dato che non esiste uno stato precedente da confrontare. La~\eqref{eq:mutex} impedisce che
un apparecchio risulti contemporaneamente acceso e spento nello stesso istante: senza di essa la
\eqref{eq:transizione} sarebbe soddisfatta anche da $u = \dw = 1$, che non ha significato fisico.

\paragraph{Durata minima di accensione.} Se il device si accende allo slot $t$, deve restare acceso
per almeno $a_i$ slot consecutivi:
\begin{equation}
  \sum_{\tau=t}^{t+a_i-1} x_{i,\tau} \;-\; a_i \, u_{i,t} \;\ge\; 0
  \qquad \forall t \in [0,\, T-a_i].
  \label{eq:durata-min}
\end{equation}
Il vincolo è attivo solo quando $u_{i,t} = 1$: in quel caso il membro sinistro impone che tutti gli
$a_i$ slot della finestra siano accesi. Quando $u_{i,t} = 0$ la disuguaglianza è soddisfatta
banalmente. Il vincolo viene generato solo se $a_i \ge 2$, perché per $a_i = 1$ sarebbe
ridondante.

\paragraph{Durata massima consecutiva.} Nessuna finestra di $b_i + 1$ slot consecutivi può essere
interamente accesa, il che equivale a limitare a $b_i$ la lunghezza di ogni ciclo:
\begin{equation}
  \sum_{\tau=t}^{t+b_i} x_{i,\tau} \;\le\; b_i
  \qquad \forall t \in [0,\, T-b_i-1].
  \label{eq:durata-max}
\end{equation}

\paragraph{Dominio.}
\begin{equation}
  x_{i,t},\, u_{i,t},\, \dw_{i,t} \in \{0,1\},
  \qquad \ep_t,\, \en_t \ge 0,
  \qquad x_{i,t} = u_{i,t} = \dw_{i,t} = 0 \;\; \forall t \notin \Vset .
\end{equation}

\section{Livelli multipli di potenza --- \code{cvxpy\_activation}}
\label{sec:m2}

\noindent\textbf{Sorgente:} \code{approach\_cvxpy\_activation.py} $\rightarrow$ \code{constrained\_highs\_multistate}\\
\textbf{Always-on:} baseline sottratta\\
\textbf{Livelli di potenza:} 3

Il modello base assume che un apparecchio assorba sempre esattamente $p_i$ watt, il che è
irrealistico. Questo modello introduce i tre livelli della~\eqref{eq:livelli}.

\subsection*{Variabile aggiuntiva e vincolo di collegamento}

Si aggiunge $z_{i,t,k} \in \{0,1\}$, che vale 1 se il device $i$ opera al livello $k$ nello slot
$t$, e si lega all'indicatore di accensione:
\begin{equation}
  x_{i,t} \;-\; \sum_{k=0}^{K-1} z_{i,t,k} \;=\; 0
  \qquad \forall i,\; \forall t .
  \label{eq:link}
\end{equation}

\begin{notebox}{Perché non serve un vincolo di mutua esclusione}
Poiché $x_{i,t}$ è binaria, la~\eqref{eq:link} impone che la somma delle $z$ valga 0 oppure 1:
\textbf{al più un livello può essere attivo}, automaticamente. Non è necessario dichiarare un
insieme SOS o aggiungere vincoli di esclusione, e questo mantiene il modello più snello.
\end{notebox}

La~\eqref{eq:link} ha una seconda funzione, altrettanto importante: $x_{i,t}$ resta l'indicatore
aggregato di accensione, quindi \textbf{tutti i vincoli di transizione e durata del
paragrafo~\ref{sec:m1} continuano a operare invariati} su di esso, senza doverli riscrivere in
termini di $z$.

\subsection*{Ricostruzione aggiornata}

L'unica altra modifica riguarda il vincolo~\eqref{eq:ricostruzione-base}, dove la potenza non è
più fissa:
\begin{equation}
  \sum_{i} \sum_{k} p_{i,k} \, z_{i,t,k} \;-\; \ep_t \;+\; \en_t \;=\; y_t
  \qquad \forall t \in \Vset .
\end{equation}

La potenza restituita in uscita è quella effettivamente stimata,
$\sum_k p_{i,k} z_{i,t,k}$, e non più il valore nominale.

\section{Apparecchi sempre accesi come variabili --- \code{cvxpy\_full}}
\label{sec:m3}

\noindent\textbf{Sorgente:} \code{approach\_cvxpy\_full.py} $\rightarrow$ \code{constrained\_highs\_full}\\
\textbf{Always-on:} modellati come variabili, nessuna baseline

Questo modello elimina il passo~1 della pipeline. Gli apparecchi si dividono in due gruppi con
trattamento diverso.

\subsection*{Apparecchi sempre accesi}

Si introduce $z^{\text{ao}}_{j,t,k}$ con il vincolo di \textbf{attività permanente}:
\begin{equation}
  \sum_{k=0}^{K-1} z^{\text{ao}}_{j,t,k} \;=\; 1
  \qquad \forall j,\; \forall t \in \Vset .
  \label{eq:always-on}
\end{equation}

Si noti che è un'\textbf{uguaglianza}, non una disuguaglianza: l'apparecchio deve trovarsi in uno
dei tre livelli in ogni istante valido, e non può spegnersi. Non ha variabili di transizione né
vincoli di durata o di fascia oraria, perché non ha cicli.

\subsection*{Apparecchi a ciclo}

Identici al paragrafo~\ref{sec:m2}: variabili $z^{\text{ev}}$, $x^{\text{ev}}$, $u^{\text{ev}}$,
$\dw^{\text{ev}}$ con l'intero corredo di vincoli.

\subsection*{Ricostruzione congiunta}

\begin{equation}
  \sum_{j} \sum_{k} p^{\text{ao}}_{j,k} \, z^{\text{ao}}_{j,t,k}
  \;+\;
  \sum_{i} \sum_{k} p_{i,k} \, z^{\text{ev}}_{i,t,k}
  \;-\; \ep_t \;+\; \en_t \;=\; y_t
  \qquad \forall t \in \Vset ,
\end{equation}
dove $y_t$ è ora il segnale \textbf{aggregato grezzo} e non più il residuo.

\begin{warnbox}{Limite strutturale}
Il vincolo~\eqref{eq:always-on} è un'uguaglianza, quindi gli apparecchi sempre accesi assorbono
per costruzione almeno $\sum_j p^{\text{ao}}_j (1-\delta)$ watt in ogni istante. Se il questionario
dichiara apparecchi la cui somma supera il consumo medio effettivo dell'abitazione, \textbf{il
modello sovrastima necessariamente} e il residuo diventa negativo.

Il caso è stato osservato su un'abitazione reale del campione: $325\unit{W}$ dichiarati
(due frigoriferi e un congelatore) contro un segnale medio di $159\unit{W}$, con un errore
sull'energia totale del $119\%$. Non è un problema di risoluzione ma di coerenza fra
questionario e misura, ed è discusso nel capitolo~\ref{ch:limiti}.
\end{warnbox}

\section{Durata flessibile --- \code{cvxpy\_soft\_duration}}
\label{sec:m4}

\noindent\textbf{Sorgente:} \code{approach\_cvxpy\_soft\_duration.py}

\subsection*{Il problema da risolvere}

Le durate dichiarate nel questionario sono approssimative. Con i vincoli rigidi
\eqref{eq:durata-min} e \eqref{eq:durata-max}, un apparecchio la cui potenza combacia esattamente
con un plateau osservato viene comunque rifiutato se il plateau dura poco più del massimo
dichiarato.

\begin{center}
\begin{minipage}{0.9\textwidth}
\itshape
Esempio reale: lavatrice con massimo dichiarato di 2 ore, plateau misurato di 2 ore e 15 minuti.
Il modello rigido copre 8 slot su 9 e lascia 15 minuti di consumo inspiegati, pur avendo a
disposizione l'apparecchio giusto.
\end{minipage}
\end{center}

\subsection*{Formulazione con variabili di scarto}

Le stesse righe di vincolo vengono rilassate aggiungendo una variabile di scarto non negativa:
\begin{align}
  \sum_{\tau=t}^{t+a_i-1} x_{i,\tau} \;+\; s^{-}_{i,t} \;-\; a_i \, u_{i,t} &\;\ge\; 0,
  & s^{-}_{i,t} &\ge 0 \label{eq:soft-min}\\
  \sum_{\tau=t}^{t+b_i} x_{i,\tau} \;-\; s^{+}_{i,t} &\;\le\; b_i,
  & s^{+}_{i,t} &\ge 0 \label{eq:soft-max}
\end{align}
con il termine di costo
\begin{equation}
  + \sum_{i} \sum_{t} \lambda_d \, p_i \, \bigl( s^{-}_{i,t} + s^{+}_{i,t} \bigr).
\end{equation}

\begin{notebox}{Il costo risulta proporzionale allo sforamento senza doverlo imporre}
Si consideri un ciclo di lunghezza $L > b_i$. Le finestre di ampiezza $b_i+1$ interamente
contenute nel ciclo sono esattamente $L - b_i$, e ciascuna richiede $s^{+} \ge 1$ per essere
soddisfatta. Il costo complessivo è quindi
\begin{equation*}
  \lambda_d \, p_i \, (L - b_i),
\end{equation*}
\textbf{lineare nello sforamento}. Simmetricamente un ciclo più corto di $a_i$ paga
$\lambda_d \, p_i \, (a_i - L)$. La proporzionalità emerge dalla struttura combinatoria dei
vincoli, non da una costruzione ad hoc.

Nell'esempio della lavatrice, uno sforamento di un solo slot costa $\lambda_d \, p_i$: abbastanza
poco da essere accettato se spiega meglio il segnale.
\end{notebox}

Il costo computazionale è trascurabile: si aggiungono due variabili \textbf{continue} per coppia
(device, slot), nessuna variabile binaria. La difficoltà del MILP resta sostanzialmente invariata.

\subsection*{Budget giornaliero (opzionale)}

Un secondo termine, disattivato per impostazione predefinita, penalizza lo scostamento del monte
ore giornaliero dall'attesa del questionario:
\begin{equation}
  E_i = \frac{\text{frequenza settimanale}}{7} \cdot \frac{\code{dur\_typical\_min}}{g},
  \qquad
  \sum_{t} x_{i,t} \;-\; d^{+}_i \;+\; d^{-}_i \;=\; E_i ,
\end{equation}
con costo $+\sum_i \lambda_D \, p_i (d^{+}_i + d^{-}_i)$.

A differenza dello scarto per ciclo questo termine è \textbf{globale}: non distingue un ciclo
lungo da molti cicli brevi. Sulle abitazioni del campione peggiorava leggermente l'errore di
ricostruzione, per cui nei modelli successivi il valore predefinito è $\lambda_D = 0$.

\begin{notebox}{Origine del limite superiore}
Nel modello rigido $b_i$ derivava da $\code{dur\_typical\_min} \times 1.5$. Qui, essendo diventato
un limite flessibile, si usa il valore \code{duration\_minutes\_max} dichiarato esplicitamente
dall'utente quando disponibile: ha senso usare il dato dichiarato ora che è possibile
oltrepassarlo pagandone il costo.
\end{notebox}

\section{Quota settimanale di utilizzi --- \code{cvxpy\_weekly\_quota}}
\label{sec:m5}

\noindent\textbf{Sorgente:} \code{approach\_cvxpy\_weekly\_quota.py}

Il questionario dichiara quante volte alla settimana ogni apparecchio viene usato. Questo modello
trasforma quel dato in un \textbf{budget scorrevole}: se la lavatrice è dichiarata due volte a
settimana e ne sono già state attribuite due, la terza accensione costa molto di più.

\subsection*{Il contatore non può essere una variabile}

Poiché ogni giorno è un modello indipendente (passo~3 della pipeline), il conteggio settimanale
\textbf{non può essere rappresentato da una variabile}: è uno stato che attraversa i solve.
La procedura è quindi ibrida.

\begin{modelbox}{Propagazione del contatore fra i giorni}
\begin{enumerate}[nosep,leftmargin=1.6em]
  \item i giorni vengono elaborati in ordine cronologico;
  \item il contatore si azzera al cambio di settimana ISO, cioè ogni lunedì;
  \item la quota residua $r_i$ entra nel modello del giorno come \textbf{costante};
  \item risolto il modello, si contano i fronti di salita nella soluzione e si aggiorna il
        contatore per il giorno successivo.
\end{enumerate}
\end{modelbox}

\subsection*{Dentro il modello: solo l'eccedenza}

Serve una sola variabile continua per apparecchio:
\begin{equation}
  \sum_{t} u_{i,t} \;-\; e_i \;\le\; r_i ,
  \qquad e_i \ge 0 ,
  \label{eq:quota}
\end{equation}
con costo $+\sum_i \lambda_{q,i} \, p_i \, e_i$.

Il meccanismo è immediato: finché il numero di accensioni resta entro $r_i$ la variabile $e_i$ vale
zero e non costa nulla; oltre quella soglia $e_i$ misura esattamente l'eccedenza. Le prime $r_i$
accensioni pagano quindi solo $\lambda_a$, quelle successive pagano anche il sovrapprezzo. Con
$r_i = 0$, cioè a quota esaurita, ogni accensione del giorno lo paga.

\subsection*{Il sovrapprezzo deve essere calibrato per apparecchio}

Il punto di pareggio di questo termine \textbf{non è 1}, a differenza degli altri. Spiegare
un'attivazione di un apparecchio che resta acceso $s_i$ slot fa guadagnare circa $s_i \cdot p_i$ di
errore di ricostruzione, quindi il sovrapprezzo deve superare quel valore per avere effetto:
\begin{equation}
  \lambda_{q,i} \;=\; \code{over\_activation\_factor} \times s_i ,
  \qquad s_i = \frac{\code{dur\_typical\_min}}{g} .
  \label{eq:quota-calibrazione}
\end{equation}

Con \code{over\_activation\_factor} $= 1.5$ la forza effettiva è la stessa per un ciclo di
microonde da 15 minuti e per uno di lavatrice da due ore. Un coefficiente scalare unico sarebbe
inevitabilmente mal calibrato: gli apparecchi del campione hanno durate tipiche che variano di un
fattore otto.

\begin{warnbox}{Allocazione greedy della quota}
Risolvendo un giorno alla volta, la quota viene spesa in modo \emph{greedy}: un giorno iniziale la
consuma liberamente e i giorni successivi della stessa settimana ne pagano le conseguenze.
Distribuirla in modo ottimo richiederebbe risolvere i sette giorni congiuntamente, con un costo
computazionale molto maggiore. La versione greedy evita comunque l'esito implausibile
``lavatrice accesa tutti i giorni'', che è lo scopo del vincolo.
\end{warnbox}

\section{Prior completi del questionario --- \code{cvxpy\_survey\_prior}}
\label{sec:m6}

\noindent\textbf{Sorgente:} \code{approach\_cvxpy\_survey\_prior.py} --- \textbf{versione in uso}

Estende il modello precedente con due termini che colmano altrettante lacune nell'uso del
questionario.

\subsection*{Stagionalità}

Il campo \code{active\_months}, che indica in quali mesi l'apparecchio viene usato, era raccolto
dal questionario ma \textbf{non veniva letto da alcun modello di ottimizzazione}. Un
climatizzatore dichiarato attivo da giugno a settembre era quindi libero di spiegare un plateau di
dicembre.

\begin{equation}
  + \sum_{i} \sum_{\substack{t \in \Vset \\ \text{mese}(t) \notin \mathcal{M}_i}}
    \lambda_s \, p_i \, x_{i,t} ,
  \label{eq:stagione}
\end{equation}
dove $\mathcal{M}_i$ è l'insieme dei mesi dichiarati. Con $\lambda_s = 8$ l'apparecchio è di fatto
rimosso fuori stagione, ma la penalità resta \textbf{flessibile}: un'evidenza schiacciante nel
segnale può ancora prevalere, coerentemente con il principio guida del documento.

\subsection*{Specificità della fascia oraria}

Un apparecchio che non aveva dichiarato alcuna fascia oraria \textbf{non pagava nulla, in nessun
momento della giornata}: era un jolly gratuito, capace di assorbire qualsiasi plateau a
qualsiasi ora. Contemporaneamente un apparecchio \emph{dentro} la propria fascia dichiarata pagava
anch'esso zero. I due pareggiavano e la scelta fra loro era arbitraria.

La correzione \textbf{non} consiste nel punire chi non ha un orario: molti apparecchi
legittimamente non ne hanno, e l'assenza di dichiarazione è assenza di informazione, non evidenza
di implausibilità. Si legge invece ``nessuna fascia'' come \textbf{``fascia di 24 ore''} e si
penalizza quanto poco la dichiarazione restringa il campo. La formulazione completa è nel
capitolo~\ref{ch:fattore}.

\subsection*{Il modello completo}

\begin{modelbox}{Funzione obiettivo della versione in uso}
\begin{align*}
  \min \;\;
  & \sum_{t \in \Vset} \bigl( \ep_t + \en_t \bigr)
  && \text{errore di ricostruzione} \\
  + \; & \sum_{i} \sum_{t \in \Vset} \lambda_w \, p_i \, f_i(t) \, x_{i,t}
  && \text{fascia oraria e specificità} \\
  + \; & \sum_{i} \sum_{t \in \Vset} \lambda_a \, p_i \, u_{i,t}
  && \text{numero di accensioni} \\
  + \; & \sum_{i} \sum_{t} \lambda_d \, p_i \bigl( s^{-}_{i,t} + s^{+}_{i,t} \bigr)
  && \text{durata dei cicli} \\
  + \; & \sum_{i} \lambda_D \, p_i \bigl( d^{+}_i + d^{-}_i \bigr)
  && \text{monte ore giornaliero} \\
  + \; & \sum_{i} \lambda_{q,i} \, p_i \, e_i
  && \text{quota settimanale} \\
  + \; & \sum_{i} \sum_{\text{fuori stagione}} \lambda_s \, p_i \, x_{i,t}
  && \text{stagionalità}
\end{align*}
\end{modelbox}

soggetta ai vincoli \eqref{eq:ricostruzione-base}--\eqref{eq:durata-max} nella forma estesa agli
apparecchi sempre accesi \eqref{eq:always-on}, con le durate rilassate
\eqref{eq:soft-min}--\eqref{eq:soft-max} e la quota settimanale \eqref{eq:quota}.

\section{Quadro riassuntivo}

\begin{center}
\begin{tabular}{@{}l c c c c c c@{}}
\toprule
\textbf{Modello} & \textbf{Always-on} & \textbf{Livelli} & \textbf{Durata} & \textbf{Quota}
& \textbf{Stagione} & \textbf{Specificità} \\
\midrule
\code{cvxpy}                & baseline  & 1 & rigida & --- & --- & --- \\
\code{cvxpy\_activation}    & baseline  & 3 & rigida & --- & --- & --- \\
\code{cvxpy\_full}          & variabili & 3 & rigida & --- & --- & --- \\
\code{cvxpy\_soft\_duration}& variabili & 3 & flessibile & --- & --- & --- \\
\code{cvxpy\_weekly\_quota} & variabili & 3 & flessibile & \checkmark & --- & --- \\
\textbf{\code{cvxpy\_survey\_prior}} & variabili & 3 & flessibile & \checkmark & \checkmark & \checkmark \\
\bottomrule
\end{tabular}
\end{center}

Le varianti con suffisso \code{\_30min}, \code{\_max} e \code{\_median} modificano soltanto la
granularità temporale o la funzione di aggregazione del ricampionamento, non la formulazione.
